\begin{lemma}
\label{lemma-algebra-map-property-in-colimit}
Let $A$ be a ring and let $B, C$ be $A$-algebras.
Suppose that $R = \colim_{i \in I} R_i$ is a directed colimit
of $A$-algebras.
\begin{enumerate}
\item If $B$ is a finite type $A$-algebra, and $u, u' : B \to C$ are
$A$-algebra maps such that
$u \otimes 1 = u' \otimes 1 : B \otimes_A R \to C \otimes_A R$
then for some $i$ we have
$u \otimes 1 = u' \otimes 1 : B \otimes_A R_i \to C \otimes_A R_i$.
\item If $C$ is a finite type $A$-algebra and $u : B \to C$ is an
$A$-algebra map such that
$u \otimes 1 : B \otimes_A R \to C \otimes_A R$ is surjective, then
for some $i$ the map $u \otimes 1 : B \otimes_A R_i \to C \otimes_A R_i$
is surjective.
\item If $C$ is of finite presentation over $A$ and
$v : C \otimes_A R \to B \otimes_A R$ is an $R$-algebra map, then there
exists an $i$ and an $R_i$-algebra map
$v_i : C \otimes_A R_i \to B \otimes_A R_i$ such that
$v = v_i \otimes 1$.
\item If $B$ is a finite type $A$-algebra, $C$ is a finitely presented
$A$-algebra, and
$u \otimes 1 : B \otimes_A R \to C \otimes_A R$ is an isomorphism, then
for some $i$ the map $u \otimes 1 : B \otimes_A R_i \to C \otimes_A R_i$
is an isomorphism.
\end{enumerate}
\end{lemma}

\begin{proof}
To prove (1) assume $u$ is as in (1) and
let $x_1, \ldots, x_m \in B$ be generators. Since
$B \otimes_A R = \colim_i B \otimes_A R_i$
we may pick an $i \in I$ such that $u(x_j) \otimes 1 = u'(x_j) \otimes 1$
in $B \otimes_A R_i$, $j = 1, \ldots, m$.
For such an $i$ we have
$u \otimes 1 = u' \otimes 1 : B \otimes_A R_i \to C \otimes_A R_i$.

\medskip\noindent
To prove (2) assume $u \otimes 1$ surjective and
let $y_1, \ldots, y_m \in C$ be generators. Since
$B \otimes_A R = \colim_i B \otimes_A R_i$
we may pick an $i \in I$ and $z_j \in B \otimes_A R_i$, $j = 1, \ldots, m$
whose images in $C \otimes_A R$ equal $y_j \otimes 1$.
For such an $i$ the map $u \otimes 1 : B \otimes_A R_i \to C \otimes_A R_i$
is surjective.

\medskip\noindent
To prove (3) let $c_1, \ldots, c_m \in C$ be generators. Let
$K = \Ker(A[x_1, \ldots, x_m] \to N)$ where the map is given by
the rule $x_j \mapsto \sum c_j$. Let $f_1, \ldots, f_t$
be generators for $K$ as an ideal in $A[x_1, \ldots, x_m]$.
We think of $f_j = f_j(x_1, \ldots, x_m)$ as a polynomial.
Since $B \otimes_A R = \colim_i B \otimes_A R_i$
we may pick an $i \in I$ and $z_j \in B \otimes_A R_i$, $j = 1, \ldots, m$
whose images in $B \otimes_A R$ equal $v(c_j \otimes 1)$.
We want to use the $z_j$ to define a map
$v_i : C \otimes_A R_i \to B \otimes_A R_i$.
Since $K \otimes_A R_i \to R_i[x_1, \ldots, x_m] \to C \otimes_A R_i \to 0$
is a presentation, it suffices to check that
$\xi_s = f_j(z_1, \ldots, z_m)$ is
zero in $B \otimes_A R_i$ for each $s = 1, \ldots, t$. This may not
be the case, but since the image of $\xi_s$ in $B \otimes_A R$ is zero
we see that it will be the case after increasing $i$ a bit.

\medskip\noindent
To prove (4) assume $u \otimes 1$ is an isomorphism, that
$B$ is a finite type $A$-algebra, and that $C$ is a finitely presented
$A$-algebra. Let $v : B \otimes_A R \to C \otimes_A R$ be an inverse to
$u \otimes 1$. Let $v_i : C \otimes_A R_i \to B \otimes_A R_i$ be as
in part (3). Apply part (1) to see that, after increasing $i$ we have
$v_i \circ (u \otimes 1) = \text{id}_{B \otimes_R R_i}$ and
$(u \otimes 1) \circ v_i = \text{id}_{C \otimes_R R_i}$.
\end{proof}
